% !TeX root = ./../../main.tex
% Rysunek własny (PHD-TS-13): przekrój kuli Blocha płaszczyzną x-z i jego obrazy
% po działaniu trzech kanałów z tabeli tab:tb:noise:channels.
% Parametry przykładowe: lambda = 0,3 (depolaryzacja), q_phi = 0,6 (tłumienie fazy),
% q_a = 0,4 (tłumienie amplitudy); półosie elips wynikają z kolumn T i t tabeli.
\begin{tikzpicture}[scale=1.25, >={Stealth[length=1.6mm]}]
  % osie
  \draw[->, gray] (-1.25,0) -- (1.35,0) node[right, black] {\footnotesize $r_x$};
  \draw[->, gray] (0,-1.25) -- (0,1.35) node[above, black] {\footnotesize $r_z$};
  % stany czyste: okrąg jednostkowy
  \draw[gray, densely dotted, line width=0.6pt] (0,0) circle (1);
  % kanał depolaryzujący: r -> (1 - lambda) r, lambda = 0,3
  \draw[line width=0.8pt] (0,0) circle (0.7);
  % tłumienie fazy: (r_x, r_z) -> (sqrt(1 - q_phi) r_x, r_z), q_phi = 0,6
  \draw[line width=0.8pt, dashed] (0,0) ellipse (0.632 and 1);
  % tłumienie amplitudy: (r_x, r_z) -> (sqrt(1 - q_a) r_x, (1 - q_a) r_z + q_a), q_a = 0,4
  \draw[line width=0.8pt, dash dot] (0,0.4) ellipse (0.775 and 0.6);
  % środki obrazów i wektor przesunięcia
  \fill (0,0) circle (0.022);
  \fill (0,0.4) circle (0.022);
  \draw[->, line width=0.6pt] (0,0) -- (0,0.378);
  \node[right] at (0.0,0.2) {\footnotesize $\ptmshift$};
  % stany bazowe
  \node[above left] at (0,1) {\footnotesize $\ket{0}$};
  \node[below left] at (0,-1) {\footnotesize $\ket{1}$};
  \node[below right] at (1,0) {\footnotesize $\ket{+}$};
  \node[below left] at (-1,0) {\footnotesize $\ket{-}$};
  % legenda po prawej stronie
  \begin{scope}[shift={(2.1,0.6)}]
    \draw[gray, densely dotted, line width=0.6pt] (0,0) -- (0.45,0) node[right, black] {\footnotesize stany czyste};
    \draw[line width=0.8pt] (0,-0.38) -- (0.45,-0.38) node[right] {\footnotesize kanał depolaryzujący, $\depolparam=0{,}3$};
    \draw[line width=0.8pt, dashed] (0,-0.76) -- (0.45,-0.76) node[right] {\footnotesize tłumienie fazy, $\qphase=0{,}6$};
    \draw[line width=0.8pt, dash dot] (0,-1.14) -- (0.45,-1.14) node[right] {\footnotesize tłumienie amplitudy, $\qamp=0{,}4$};
  \end{scope}
\end{tikzpicture}
